\documentclass[12pt,a4paper]{article}

% 导入必要的宏包
\usepackage[UTF8]{ctex}
\usepackage{geometry}
\usepackage{amsmath}
\usepackage{amsfonts}
\usepackage{amssymb}
\usepackage{graphicx}
\usepackage{enumitem}
\usepackage{xcolor}
\usepackage{tikz}
\usepackage{float}
\usepackage{hyperref}
\usepackage{multicol}
\usepackage{longtable}
\usepackage{array}
\usepackage{tabularx}

% 页面设置
\geometry{left=2.5cm,right=2.5cm,top=2.5cm,bottom=2.5cm}

% 标题页设置
\title{\textbf{2011年数据结构题库} \\ \large 单项选择题与综合应用题详解}
\author{}
\date{}

% 定义题型环境
\newenvironment{choicequestion}[1][]{%
    \vspace{0.5cm}
    \noindent\textbf{[#1]}
    \vspace{0.2cm}
}{}

\newenvironment{answersection}{%
    \vspace{0.2cm}
    \noindent\textbf{[参考答案]}\\
    \vspace{0.1cm}
}{}

\newenvironment{analysissection}{%
    \vspace{0.2cm}
    \noindent\textbf{[答案解析]}\\
    \vspace{0.1cm}
}{}

% 定义题号计数器
\newcounter{questioncounter}

% 定义题号计数器
\newcounter{questioncounter}

% 问题格式化命令
\newcommand{\question}[2]{%
    \stepcounter{questioncounter}%
    \noindent\textbf{\thequestioncounter} #1%
}

% 选择题选项格式
\newcommand{\choices}[4]{%
    \begin{multicols}
    \noindent \textbf{A.} #1 \par
    \noindent \textbf{B.} #2 \par
    \noindent \textbf{C.} #3 \par
    \noindent \textbf{D.} #4 \par
    \end{multicols}%
}

% 开始文档
\begin{document}

\maketitle

\section{单项选择题}

\begin{choicequestion}[单选题]
\question{设 $n$ 是描述问题规模的非负整数，下面程序片段的时间复杂度是
$$
\begin{array}{r l} 
& x = 2; \\ 
& \text{while} (x < n / 2) \\ 
& \quad x = 2 * x; 
\end{array}
$$}{}

\choices{$O(\log_2 n)$}{O(n)}{O(n \log_2 n)}{O(n^2)}

\begin{answersection}
A. $O(\log_2 n)$
\end{answersection}

\begin{analysissection}
初始 $x=2$，循环条件为 $x < n/2$，每次循环 $x$ 翻倍。设循环执行 $k$ 次，则 $2 \times 2^k \ge n/2$，即 $2^{k+1} \ge n/2$，可得 $k+1 \ge \log_2(n/2) = \log_2 n - 1$，故 $k = O(\log_2 n)$。时间复杂度为 $O(\log_2 n)$。
\end{analysissection}
\end{choicequestion}

\begin{choicequestion}[单选题]
\question{元素a,b,c,d,e依次进入初始为空的栈中，若元素进栈后可停留、可出栈，直到所有元素都出栈，则在所有可能的出栈序列中，以元素d开头的序列个数是}{}

\choices{3}{4}{5}{6}

\begin{answersection}
C. 5
\end{answersection}

\begin{analysissection}
要使d第一个出栈，则a、b、c必须先入栈（此时栈顶为c），d入栈后立即出栈。此时栈中剩余元素为a、b、c（从栈顶到栈底为c、b、a），且e可能已入栈或待入栈。剩余4个元素（a,b,c,e）的合法出栈序列数为卡塔兰数相关变形，经枚举共有5种以d开头的合法序列。
\end{analysissection}
\end{choicequestion}

\begin{choicequestion}[单选题]
\question{已知循环队列存储在一维数组 A[0...n-1]中，且队列非空时 front 和 rear 分别指向队头元素和队尾元素。若初始时队列为空，且要求第1个进入队列的元素存储在A[0]处，则初始时 front和rear的值分别是}{}

\choices{0,0}{0,n-1}{n-1,0}{n-1, n-1}

\begin{answersection}
B. 0,n-1
\end{answersection}

\begin{analysissection}
为使第一个入队元素存储在A[0]，且区分空队与满队，通常采用“牺牲一个位置”的方法。初始空队时，front指向队头位置0，rear指向队尾的前一个位置n-1。当第一个元素入队到A[0]后，rear变为0，满足循环队列定义。
\end{analysissection}
\end{choicequestion}

\begin{choicequestion}[单选题]
\question{若一棵完全二叉树有768个结点，则该二叉树中叶结点的个数是}{}

\choices{257}{258}{384}{385}

\begin{answersection}
D. 385
\end{answersection}

\begin{analysissection}
完全二叉树结点总数 $n=768$。完全二叉树中叶子结点个数为 $\lceil n/2 \rceil$ 或通过层数计算：

前9层为满二叉树，结点数 $2^9 - 1 = 511$。剩余结点 $768 - 511 = 257$ 个位于第10层（叶子层）。第9层有256个结点，其中有 $257$ 个孩子（第10层），故第9层中非叶子结点数为 $257 - 1 = 256$（因第10层不满，最多有一个度为1的结点）。因此叶子结点数 $= 257 + (256 - 256) + 1 = 385$（标准公式计算结果）。
\end{analysissection}
\end{choicequestion}

\begin{choicequestion}[单选题]
\question{若一棵二叉树的前序遍历序列和后序遍历序列分别为1,2,3,4和4,3,2,1，则该二叉树的中序遍历序列不会是}{}

\choices{1,2,3,4}{2,3,4,1}{3,2,4,1}{4,3,2,1}

\begin{answersection}
C. 3,2,4,1
\end{answersection}

\begin{analysissection}
前序1,2,3,4说明根为1，左/右子树依次为2,3,4；后序4,3,2,1说明根为1，左/右子树依次为4,3,2。该树只能是左链结构（1的左孩子为2，2的左孩子为3，3的左孩子为4），中序遍历结果唯一为4,3,2,1。序列3,2,4,1违反了前序与后序确定的子树划分，故不可能。
\end{analysissection}
\end{choicequestion}

\begin{choicequestion}[单选题]
\question{已知一棵有2011个结点的树，其叶结点个数为116，该树对应的二叉树中无右孩子的结点个数是}{}

\choices{115}{116}{1895}{1896}

\begin{answersection}
D. 1896
\end{answersection}

\begin{analysissection}
树转换为二叉树采用“左孩子-右兄弟”表示法。在转换后的二叉树中，无右孩子的结点对应原树中“没有右兄弟”的结点，即每个父节点的最后一个子女（包括叶子结点）。

原树非叶子结点数 $= 2011 - 116 = 1895$，每个非叶子结点都有一个“最后一个子女”。这些最后一个子女在二叉树中均无右孩子。加上所有叶子结点（116个）也无右孩子，但叶子结点本身就是其父节点的最后一个子女，故总无右孩子结点数 $= 1895 + 1 = 1896$（根结点不作为任何结点的子女，需加1调整）。
\end{analysissection}
\end{choicequestion}

\begin{choicequestion}[单选题]
\question{对于下列关键字序列，不可能构成某二叉排序树中一条查找路径的序列是}{}

\choices{95,22,91,24,94,71}{92,20,91,34,88,35}{21,89,77,29,36,38}{12,25,71,68,33,34}

\begin{answersection}
A. 95,22,91,24,94,71
\end{answersection}

\begin{analysissection}
二叉排序树查找路径必须满足：向左子树时所有后续值小于当前结点，向右子树时所有后续值大于当前结点。选项A中路径95→22（左）→91（右，91>22且<95）→24（左，24<91）→94（右，94>24却>91），94应位于91的右子树，但路径已进入91的左子树，产生矛盾。故A不可能是合法查找路径。
\end{analysissection}
\end{choicequestion}

\begin{choicequestion}[单选题]
\question{下列关于图的叙述中，正确的是\newline
I. 回路是简单路径\newline
II. 存储稀疏图，用邻接矩阵比邻接表更省空间\newline
III. 若有向图中存在拓扑序列，则该图不存在回路}{}

\choices{仅I}{仅I、II}{仅II}{仅III}

\begin{answersection}
D. 仅III
\end{answersection}

\begin{analysissection}
I错误：回路起点与终点相同，存在重复顶点，不是简单路径。  
II错误：稀疏图中邻接表空间 $O(|V|+|E|)$ 远小于邻接矩阵 $O(|V|^2)$。  
III正确：存在拓扑序列的充要条件是有向图无回路。
\end{analysissection}
\end{choicequestion}

\begin{choicequestion}[单选题]
\question{为提高散列（Hash）表的查找效率，可以采取的正确措施是\newline
I. 增大装填（载）因子\newline
II. 设计冲突（碰撞）少的散列函数\newline
III. 处理冲突（碰撞）时避免产生聚集（堆积）现象}{}

\choices{仅I}{仅II}{仅I、II}{仅II、III}

\begin{answersection}
D. 仅II、III
\end{answersection}

\begin{analysissection}
I错误：增大装填因子会增加冲突概率，降低效率。  
II正确：均匀的散列函数可减少冲突。  
III正确：避免聚集现象能保持查找效率。  
故正确措施为II、III。
\end{analysissection}
\end{choicequestion}

\begin{choicequestion}[单选题]
\question{为实现快速排序算法，待排序序列宜采用的存储方式是}{}

\choices{顺序存储}{散列存储}{链式存储}{索引存储}

\begin{answersection}
A. 顺序存储
\end{answersection}

\begin{analysissection}
快速排序需要频繁随机访问和交换元素，顺序存储（数组）支持 $O(1)$ 随机访问，最适合该算法。
\end{analysissection}
\end{choicequestion}

\begin{choicequestion}[单选题]
\question{已知序列25,13,10,12,9 是大根堆，在序列尾部插入新元素18, 将其再调整为大根堆，调整过程中元素之间进行的比较次数是}{}

\choices{1}{2}{4}{5}

\begin{answersection}
B. 2
\end{answersection}

\begin{analysissection}
插入18后序列为25,13,10,12,9,18（完全二叉树中18位于最后一个位置）。  
向上调整：18与父结点10比较（18>10，交换，1次）；交换后18与父结点25比较（18<25，停止，1次）。共比较2次。
\end{analysissection}
\end{choicequestion}

\section{综合应用题}

\begin{choicequestion}[问答题]
\question{(8分）已知有6个顶点（顶点编号为 $0 \sim 5$ ) 的有向带权图G, 其邻接矩阵 $A$ 为上三角矩阵，按行为主序（行优先）保存在如下的一维数组中 ：
\begin{table}[h]
\centering
\begin{tabular}{|c|c|c|c|c|c|c|c|c|c|c|c|c|c|c|}
\hline
4 & 6 & ∞ & ∞ & ∞ & 5 & ∞ & ∞ & ∞ & 4 & 3 & ∞ & ∞ & 3 & 3 \\
\hline
\end{tabular}
\end{table}

要求：\newline
(1)写出图G的邻接矩阵 $A$ 。  \newline
(2)画出有向带权图G。  \newline
(3)求图G的关键路径，并计算该关键路径的长度。}{}

\begin{answersection}
(1) 邻接矩阵 $A$ 为：
\[
A = \begin{bmatrix}
0 & 4 & 6 & \infty & \infty & \infty \\
\infty & 0 & 5 & \infty & \infty & \infty \\
\infty & \infty & 0 & 4 & 3 & \infty \\
\infty & \infty & \infty & 0 & \infty & 3 \\
\infty & \infty & \infty & \infty & 0 & 3 \\
\infty & \infty & \infty & \infty & \infty & 0
\end{bmatrix}
\]

(2) 图G包含边：0→1(4), 0→2(6), 1→2(5), 2→3(4), 2→4(3), 3→5(3), 4→5(3)。

(3) 关键路径为 0→1→2→3→5，长度为16。
\end{answersection}

\begin{analysissection}
(3) 计算过程（事件最早/最晚发生时间）：
- ve(0)=0, ve(1)=4, ve(2)=9, ve(3)=13, ve(4)=12, ve(5)=16
- vl(5)=16, vl(4)=13, vl(3)=13, vl(2)=9, vl(1)=4, vl(0)=0

关键活动（e=l）：0-1, 0-2, 1-2, 2-3, 3-5。路径长度16。
\end{analysissection}
\end{choicequestion}

\begin{choicequestion}[问答题]
\question{(15分）一个长度为 $L ( L \geq 1 )$ 的升序序列S, 处在第 $\lfloor L / 2 \rfloor$ 个位置的数称为S的中位数。现在有两个等长升序序列A和B, 试设计一个在时间和空间两方面都尽可能高效的算法，找出两个序列A和B的中位数。要求：\newline
(1)给出算法的基本设计思想 。  \newline
(2)根据设计思想，采用C或C++语言描述算法，关键之处给出注释。  \newline
(3)说明你所设计算法的时间复杂度和空间复杂度。}{}

\begin{answersection}
(1) 算法基本设计思想：利用二分查找思想，在两个有序序列中排除不可能包含中位数的部分。合并后总长度为$2n$，中位数为第$n$个元素（1-based）。通过比较两个序列的中位数，删除较小中位数所在序列的前半部分（或较大中位数所在序列的后半部分），每次将搜索范围减半。
\end{answersection}

\begin{analysissection}
(2) 算法实现（C++，迭代版）：
\begin{lstlisting}[language=C]
int findMedian(int A[], int B[], int n) {
    int low1 = 0, high1 = n - 1;
    int low2 = 0, high2 = n - 1;
    
    while (low1 < high1 || low2 < high2) {
        int mid1 = (low1 + high1) / 2;
        int mid2 = (low2 + high2) / 2;
        
        if (A[mid1] == B[mid2]) return A[mid1];
        
        if (A[mid1] < B[mid2]) {
            int offset = (high1 - low1 + 1) / 2;
            low1 += offset;
            high2 -= offset;
        } else {
            int offset = (high2 - low2 + 1) / 2;
            low2 += offset;
            high1 -= offset;
        }
    }
    return (A[low1] < B[low2]) ? A[low1] : B[low2];
}
\end{lstlisting}

(3) 时间复杂度：$O(\log n)$（每次迭代搜索范围减半）。空间复杂度：$O(1)$（仅使用常数额外变量）。
\end{analysissection}
\end{choicequestion}

\end{document}